% Elevation sketch of the contribution the normal component of the commanded
% force makes through a tangential compliance-centre displacement. The view is
% taken in the surface tangent plane, with n_s up, so the two panels differ only
% in the direction of r_{c,t}: the tool angular offset and force are identical, and
% the moment contribution therefore reverses with the displacement.
% Uses only tikz + arrows.meta/calc, which config/packages.tex loads.
%
% Colour follows the drawn palette and matches the moment-bookkeeping figure:
% blue for the virtual displacement r_{c,t}, red for a force, black for the
% moment and the axes. No line is grey or white, and every line is solid.
%
% Layout note. Each panel is kept to about 5.2 cm wide and 2.7 cm tall so the
% pair sits near \textwidth without being scaled down; an earlier wider version
% was shrunk by \resizebox until its labels collided. Three heights are used and
% nothing else may be placed on them: the displacement at the TCP height, the
% force tail above it, and the moment arc above that.
\begin{tikzpicture}[
    font=\small,
    surf/.style={draw=black, fill=black!8},
    tool/.style={draw=black, fill=black!4},
    frc/.style={-{Latex[length=2.2mm]}, thick, red},
    lev/.style={-{Latex[length=2.2mm]}, thick, blue},
    mom/.style={-{Latex[length=2.2mm]}, thick, black},
    ax/.style={-{Latex[length=1.8mm]}, black},
    lbl/.style={font=\footnotesize, inner sep=1.8pt},
    pt/.style={circle, fill=black, inner sep=0pt, minimum size=3.4pt},
    cpt/.style={circle, fill=blue, inner sep=0pt, minimum size=3.4pt},
  ]

% ============ (a) displacement selected from the offset direction ===========
\begin{scope}

  % ---- surface plane ----------------------------------------------------
  \fill[surf] (-2.6,-0.34) rectangle (2.6,0);
  \draw[black] (-2.6,0) -- (2.6,0);

  % ---- surface normal ----------------------------------------------------
  \draw[ax] (-2.45,0) -- (-2.45,0.80) node[lbl, above] {$n_s$};

  % ---- tool face, inclined and contacting at its lower edge --------------
  \begin{scope}[shift={(0,0.125)}, rotate=-10]
    \draw[tool] (-0.72,0) rectangle (0.72,0.30);
    \coordinate (T) at (0,0.30);
  \end{scope}

  % ---- compliance centre and its tangential displacement, at TCP height --
  \coordinate (P) at ($(T)+(-2.0,0)$);
  \draw[lev] (T) -- (P);
  \node[lbl, above, blue] at ($(P)!0.32!(T)+(0,0.04)$) {$r_{c,t}$};
  \node[cpt] at (P) {};
  \node[lbl, left, blue] at (P) {$p_c$};
  \node[pt] at (T) {};
  \node[lbl, anchor=south east] at ($(T)+(-0.07,0.05)$) {TCP};

  % ---- normal component of the commanded force, pressing into the plane --
  \draw[frc] ($(T)+(0,1.05)$) -- ($(T)+(0,0.14)$);
  \node[lbl, right, red] at ($(T)+(0.72,0.60)$) {$f_n$};

  % ---- the contribution it makes through the displacement ----------------
  \coordinate (M) at ($(T)+(0,0.30)$);
  \draw[mom] ([shift=(30:0.95)]M) arc (30:150:0.95);
  \node[lbl] at ($(T)+(0,1.80)$) {$r_{c,t}\times f_n$};

  \node at (0,-0.92) {(a) Selected displacement};
\end{scope}

% ============ (b) displacement in the opposite direction ====================
\begin{scope}[xshift=6.0cm]

  \fill[surf] (-2.6,-0.34) rectangle (2.6,0);
  \draw[black] (-2.6,0) -- (2.6,0);

  \draw[ax] (-2.45,0) -- (-2.45,0.80) node[lbl, above] {$n_s$};

  \begin{scope}[shift={(0,0.125)}, rotate=-10]
    \draw[tool] (-0.72,0) rectangle (0.72,0.30);
    \coordinate (T2) at (0,0.30);
  \end{scope}

  % ---- the displacement is reversed; everything else is unchanged --------
  \coordinate (P2) at ($(T2)+(2.0,0)$);
  \draw[lev] (T2) -- (P2);
  \node[lbl, above, blue] at ($(P2)!0.32!(T2)+(0,0.04)$) {$r_{c,t}$};
  \node[cpt] at (P2) {};
  \node[lbl, right, blue] at (P2) {$p_c$};
  \node[pt] at (T2) {};
  \node[lbl, anchor=south east] at ($(T2)+(-0.07,0.05)$) {TCP};

  \draw[frc] ($(T2)+(0,1.05)$) -- ($(T2)+(0,0.14)$);
  \node[lbl, right, red] at ($(T2)+(0.72,0.60)$) {$f_n$};

  \coordinate (M2) at ($(T2)+(0,0.30)$);
  \draw[mom] ([shift=(150:0.95)]M2) arc (150:30:0.95);
  \node[lbl] at ($(T2)+(0,1.80)$) {$r_{c,t}\times f_n$};

  \node at (0,-0.92) {(b) Reversed displacement};
\end{scope}

\end{tikzpicture}
